% !TeX root = ../wangshuo_phdthesis.tex

\begin{acknowledgements}

  在北京师范大学求学的这些年，是我人生中一段意义深远的旅程。回望求索之路的点滴积累，心中满怀感恩。

  衷心感谢我的导师张江教授。张老师在论文写作与研究推进中倾注了大量心血，从科学问题的凝练、创新逻辑的构建，到每一处表述的精准推敲，始终以极高的标准要求我反复打磨、持续迭代。正是在这样严格而耐心的训练中，我完成了从工程思维向科研思维的转变——不仅学会了解决问题，更学会了提出有价值的问题，将零散的技术突破提升为系统性的认知。张老师让我理解，真正的学术追求源于对世界运行机制的不懈追问，以及对认知边界的不断逼近。

  感谢国家留学基金委（CSC）的资助，使我有机会赴苏黎世联邦理工学院开展联合培养研究。Lothar Thiele教授与程云老师的悉心指导，开阔了我的国际学术视野。感谢樊京芳教授、张钢锋、滕梦凡、谷伟伟等老师在人工智能与地球系统科学交叉方向上的合作，加深了我对学科融合的认识。感谢新加坡国立大学杨柳教授，他对问题本质的深刻洞察以及在ICON等前沿技术上的探索，为我提供了重要的启发。

  感谢张江老师与张倩搭建的集智俱乐部学术交流平台，这个跨学科的开放社区为我提供了丰富的思想碰撞空间，也让我结识了许多志同道合的朋友。感谢预答辩专家小组的老师提出的宝贵意见。感谢许菁、高飞、刘晶、陈昊、王志鹏、张章、任汉承等同学在合作研究中的支持与启发。感谢彩云科技的同事们，共事的经历不仅锻炼了我的工程实践能力，也让我在服务千万级用户的过程中切实体会到科研成果落地的价值。

  我始终相信，论文应当写在祖国大地上。当看到自己研发的模型被应用于水利管理与大气污染监测部门时，那种满足感远胜于论文发表本身。从参与水利部门的技术交流，到关注观测资源匮乏地区的环境问题，推动环境公平、服务公共健康，是我持续前行的核心动力。我将坚定走产学研融合之路，坚持做算法有创新、成果能落地的研究。

  荣格曾言，完整大于完美。回望读博历程，困顿与挑战让我更为坚韧，也让我更加完整。在苏黎世机场，我曾看到一句话：\emph{Think Long-term, Create Long-term Value}，这与我内心的体悟深深呼应。找到属于自己的Ikigai——去做热爱的、擅长的、世界真正需要的事情——是我在这段旅程中逐渐确立的信念。读博如修行，亦如人生。

  感谢我的爱人秦颖女士一路以来的理解、支持与陪伴。感谢我的孩子王行知——"行知"承载着我的人生信念：知行合一，行是知之始。愿孩子无论面对何种困难，都能保持探索的勇气与内心的良知。

  感谢所有在我生命中出现过的人与事。

\end{acknowledgements}
